\documentclass[11pt]{article}

% ---------- Preámbulo ----------
\usepackage[utf8]{inputenc}
\usepackage[T1]{fontenc}
\usepackage[spanish]{babel}
\usepackage[margin=2.5cm]{geometry}
\usepackage{amsmath,amssymb}
\usepackage{siunitx}
\usepackage{booktabs}
\usepackage{graphicx}
\usepackage[backend=biber,style=ieee,sorting=none]{biblatex}
\usepackage{hyperref}
\addbibresource{references.bib}

\title{TA-IND-04 --- Informe Técnico Individual\\
\large Análisis de Rendimiento Paralelo aplicado al Proyecto Fin de Curso}
\author{}
\date{}

\begin{document}

% ============================================================
% PORTADA
% ============================================================
\begin{titlepage}
\centering
{\large UNIVERSIDAD TÉCNICA ESTATAL DE QUEVEDO}\\
Facultad de Ciencias de la Computación\\
Carrera de Ingeniería de Software\\[1.5em]

{\Large\bf TA-IND-04 --- Informe Técnico Individual}\\[0.3em]
{\large Análisis de Rendimiento Paralelo (Unidad 4)\\
aplicado al Proyecto Fin de Curso}\\[1.5em]

Asignatura: Aplicaciones Distribuidas (ISR-701)\\
Unidad 4 --- Cómputo Paralelo y Distribuido\\
Período académico: 2026--2027 PPA\\[2em]

\begin{tabular}{ll}
\textbf{Estudiante:} & Cristhian Daniel Pacheco Cárdenas \\
\textbf{Equipo de PE-U4:} & Equipo A --- Cristhian Daniel Pacheco Cárdenas, \\
 & Robinson Rodrigo Cando Moreno, Ernesto Gregory Luna Mora \\
\textbf{PFC de referencia:} & ACC --- Sistema de Gestión de Soporte Técnico ISP \\
\textbf{Transformación foco:} & T3 (Join, \emph{tickets} $\bowtie$ \emph{agentes}) \\
\textbf{Repositorio individual:} & \url{https://github.com/CristhianP03/ta-ind-04-pacheco} \\
\textbf{Fecha:} & 8 de agosto de 2026 \\
\end{tabular}

\vfill
\textbf{Docente:} Gleiston C. Guerrero-Ulloa, M.Sc.
\end{titlepage}

\section{Introducción y contexto}

Este informe analiza, de forma individual, los resultados de speedup
obtenidos en la práctica grupal GA-SUM-05/PE-U4, en la que el equipo A
implementó y midió cinco transformaciones (T1--T5) sobre un dataset
sintético de 600\,000 tickets de soporte técnico ISP, comparando pandas
(línea base de una sola máquina) contra PySpark 3.5.3 en Google Colab,
con escalado controlado de T3 (unión \emph{tickets}$\bowtie$\emph{agentes})
en $N=1,2,4$ executors. El objetivo de este documento no es repetir ese
experimento grupal, sino someter sus resultados a un análisis
cuantitativo propio frente al modelo de Amdahl, diagnosticar la fracción
no escalable observada, y fundamentar --- con datos propios de PE-U4 y
la arquitectura real del Proyecto Fin de Curso --- una propuesta
concreta de integración del procesamiento distribuido en dicho proyecto.

Los datos base se tomaron del repositorio
\url{https://github.com/CristhianP03/pe-u4-spark-equipoA}, commit
\texttt{aac987686e1d8dd0559b3ea6ed07af3335570ced}, archivo
\texttt{resultados/tiempos\_resumen.csv}, verificado línea por línea
contra las mediciones crudas del mismo repositorio antes de su uso en
este documento. Se descartó explícitamente el uso de cualquier tabla o
cálculo del informe grupal (\texttt{docs/PE\_U4\_Informe.pdf}) más allá
de los datos crudos de tiempos: el ajuste de Amdahl, la métrica de
Karp-Flatt y el umbral de rentabilidad presentados aquí emplean una
metodología propia, distinta de la reportada por el equipo, como se
declara explícitamente en la Sección~9.

El dominio de referencia es \textbf{ACC --- Sistema de Gestión de
Soporte Técnico ISP}, un sistema de microservicios (Spring Boot, Node.js
y FastAPI) con persistencia sobre un clúster CockroachDB de 3 nodos, en
cuya Entrega 3 ya existe un \emph{pipeline} de Spark propio para
detección de reincidencia de clientes y \emph{clustering} de tickets ---
pipeline cuya integración formal, periodicidad, dimensionamiento y
consumidor de salida son precisamente el objeto de la Sección~5 de este
informe.

\section{Análisis propio de los resultados de speedup}

\subsection{Marco matemático}

La ley de Amdahl relaciona el speedup $S(N)$ obtenido con $N$ unidades
de procesamiento y la fracción paralelizable $p$ del cómputo:

\begin{equation}
S(N) = \frac{1}{(1-p) + \dfrac{p}{N}}, \qquad
S_{\max} = \lim_{N\to\infty} S(N) = \frac{1}{1-p}.
\label{eq:amdahl}
\end{equation}

\subsection{Datos base}

Los tiempos de ejecución se tomaron del commit
\texttt{aac987686e1d8dd0559b3ea6ed07af3335570ced} del repositorio de
PE-U4 del equipo A, archivo \texttt{resultados/tiempos\_resumen.csv},
generado sobre Google Colab con PySpark 3.5.3. Cada valor corresponde a
la mediana de 5 repeticiones cronometradas, tras descartar una ejecución
de calentamiento por combinación (transformación, motor).

\begin{table}[h]
\centering
\caption{Tiempos medianos y speedup medido por transformación.}
\label{tab:tiempos}
\begin{tabular}{lccS[table-format=2.3]S[table-format=1.3]}
\toprule
Transf. & Motor & $N$ & {$T$ (\si{\second})} & {$S(N)$} \\
\midrule
T1 & pandas & --- & 9.437  & --- \\
T1 & Spark  & 4    & 6.210  & 1.520 \\
T2 & pandas & --- & 8.345  & --- \\
T2 & Spark  & 4    & 4.363  & 1.913 \\
T3 & pandas & --- & 10.119 & --- \\
T3 & Spark  & 1    & 7.891  & 1.282 \\
T3 & Spark  & 2    & 7.059  & 1.433 \\
T3 & Spark  & 4    & 7.087  & 1.428 \\
T4 & pandas & --- & 12.296 & --- \\
T4 & Spark  & 4    & 8.381  & 1.467 \\
T5 & pandas & --- & 6.269  & --- \\
T5 & Spark  & 4    & 2.450  & 2.558 \\
\bottomrule
\end{tabular}
\end{table}

\subsection{Análisis agregado de las cinco transformaciones}

Las cinco transformaciones muestran speedups a $N=4$ entre 1.467 (T4,
columna derivada) y 2.558 (T5, top-\emph{N}). T5 obtiene el mayor
beneficio porque su costo dominante es una agregación con muy pocas
filas de salida, que Spark resuelve con un \emph{shuffle} reducido tras
la fase de agregación parcial por partición. T4, en cambio, aplica
transformaciones fila por fila sin reordenar datos entre particiones;
pandas ya ejecuta esa operación de forma vectorizada y eficiente en
memoria, de modo que el margen de mejora al distribuirla es menor.

\subsection{Desarrollo matemático completo para la transformación foco (T3)}

Para T3 se dispone de mediciones con $N=1,2,4$ executors, lo que permite
ajustar la fracción paralelizable $p$ de la ecuación~\eqref{eq:amdahl}
por mínimos cuadrados. Se excluye deliberadamente el punto $N=1$ del
ajuste: la ley de Amdahl impone $S(1)=1$ por construcción (sin
paralelismo no hay aceleración posible), mientras que el valor medido es
$S(1)=1.282$; esa diferencia no proviene de paralelismo --- imposible
con un único executor --- sino de que Spark y pandas son motores de
ejecución distintos, y el optimizador de consultas de Spark ya introduce
ganancias de motor independientes del número de executors. Incluir ese
punto en el ajuste distorsionaría la estimación de $p$ sin aportar
información sobre escalabilidad real.

Minimizando $\sum_{N\in\{2,4\}}\left(S_{\mathrm{med}}(N)-S(N;p)\right)^2$
se obtiene $\hat{p}=0.4445$, con lo que la ecuación~\eqref{eq:amdahl}
predice $S_{\mathrm{teo}}(2)=1.286$, $S_{\mathrm{teo}}(4)=1.500$ y un
techo asintótico $S_{\max}=1/(1-\hat{p})=1.800$.

\begin{table}[h]
\centering
\caption{Speedup medido frente al predicho por la ley de Amdahl ($\hat p = 0.4445$).}
\label{tab:amdahl}
\begin{tabular}{cS[table-format=1.3]S[table-format=1.3]S[table-format=-2.2]}
\toprule
$N$ & {$S_{\mathrm{med}}$} & {$S_{\mathrm{teo}}$} & {$\Delta$ (\%)} \\
\midrule
2 & 1.433 & 1.286 & +11.49 \\
4 & 1.428 & 1.500 & -4.82  \\
\bottomrule
\end{tabular}
\end{table}

\begin{figure}[h]
\centering
\includegraphics[width=0.85\linewidth]{figuras/fig_speedup.png}
\caption{Speedup medido de T3 (unión pandas/Spark) frente a la curva
teórica de Amdahl ajustada por mínimos cuadrados ($\hat{p}=0.4445$).
La brecha entre la curva de paralelismo ideal ($S=N$) y la curva teórica
refleja el límite estructural impuesto por la fracción no paralelizable;
la brecha entre la curva teórica y los puntos medidos refleja la
sobrecarga real del motor, que se cuantifica en la Sección~4 mediante
la métrica de Karp-Flatt.}
\label{fig:speedup_t3}
\end{figure}

\subsection{Respuesta a la pregunta 1 de la guía}

Los resultados no coinciden con la predicción de Amdahl de forma
uniforme: a $N=2$ el speedup medido \emph{supera} al teórico en un
$11.49\,\%$, mientras que a $N=4$ queda \emph{por debajo} en un
$4.82\,\%$ (Tabla~\ref{tab:amdahl}, Figura~\ref{fig:speedup_t3}). Esta
inversión de signo es, en sí misma, el hallazgo relevante. Con dos
executors, la ganancia adicional de paralelismo todavía compensa la
sobrecarga de coordinación, y el optimizador de Spark aporta una mejora
de motor no capturada por el modelo de Amdahl. Con cuatro executors, en
cambio, el costo de coordinar más particiones --- reparto de trabajo
del \emph{join}, transferencia de datos entre executors y planificación
de tareas --- empieza a superar la ganancia marginal de paralelismo,
algo que Amdahl no puede anticipar porque asume implícitamente que la
única variable es $N$ y que el costo de coordinación es despreciable
\cite{amdahl1967}. Hill y Marty señalan precisamente que este tipo de
rendimientos decrecientes se agrava cuanto mayor es el número de
unidades coordinadas, incluso en arquitecturas de propósito general
\cite{hillmarty2008}. La Sección~4 retoma esta misma inversión con la
métrica de Karp-Flatt, que confirma cuantitativamente que la causa es la
sobrecarga creciente del motor y no un límite algorítmico fijo.


\section{Diagnóstico de la fracción no escalable}


\subsection{Métrica de Karp-Flatt}

La fracción serial experimental $e$, determinada a partir del speedup
observado $S$ con $N$ unidades, se define como:

\begin{equation}
e = \frac{\dfrac{1}{S} - \dfrac{1}{N}}{1 - \dfrac{1}{N}}, \qquad N > 1.
\label{eq:karpflatt}
\end{equation}

Sustituyendo en la ecuación~\eqref{eq:karpflatt} los valores medidos de
T3 (Tabla~\ref{tab:amdahl}):

\begin{align}
e(2) &= \frac{\frac{1}{1.43346} - \frac{1}{2}}{1 - \frac{1}{2}}
      = \frac{0.69761 - 0.5}{0.5} = 0.39522 \\
e(4) &= \frac{\frac{1}{1.42772} - \frac{1}{4}}{1 - \frac{1}{4}}
      = \frac{0.70042 - 0.25}{0.75} = 0.60055
\end{align}

La fracción serial experimental \textbf{crece un 51.95\,\%} al pasar de
$N=2$ a $N=4$ ($0.395\rightarrow0.601$). Según el criterio de
interpretación de Karp y Flatt \cite{karpflatt1990}, un valor de $e$
que se mantiene aproximadamente constante señala un límite algorítmico
fijo, mientras que un valor creciente indica sobrecarga adicional del
motor de ejecución al aumentar $N$. El fuerte incremento observado
descarta la primera hipótesis: la pérdida de escalabilidad de T3 entre
$N=2$ y $N=4$ no es atribuible a la fracción secuencial intrínseca del
algoritmo de \emph{join}, sino a costos de coordinación que crecen con
el número de executors.

\subsection{Eficiencia del paralelismo}

La eficiencia mide qué fracción del beneficio ideal se aprovecha
realmente por cada unidad de procesamiento adicional:

\begin{equation}
E(N) = \frac{S(N)}{N}.
\label{eq:eficiencia}
\end{equation}

Con los valores medidos de T3, $E(1)=1.282$, $E(2)=0.717$ y
$E(4)=0.357$ (ecuación~\eqref{eq:eficiencia}): la eficiencia cae de
forma pronunciada y monótona incluso donde el speedup absoluto todavía
crece (de $N=1$ a $N=2$), lo que confirma visualmente el mismo
fenómeno que Karp-Flatt cuantifica --- cada executor adicional aporta
una fracción decreciente de beneficio real.

\subsection{Atribución a mecanismos concretos del motor}

La captura de la Spark UI para T3 con $N=4$ (\texttt{evidencia/spark\_ui\_t3.png}
del repositorio de PE-U4, \emph{query} SQL de duración total 12\,s)
permite identificar cuatro mecanismos concretos que explican el
crecimiento de $e$:

\begin{enumerate}
  \item \textbf{Colapso adaptativo de particiones (AQE).} El nodo
  \texttt{AQEShuffleRead} reporta \texttt{number of coalesced
  partitions: 1} en ambos lados del \emph{join}. Aunque la sesión de
  Spark se configuró con 4 executors, el motor de ejecución adaptativa
  reunió el \emph{shuffle} completo en una única partición por el
  volumen reducido de datos (180\,003 registros tras el filtrado), de
  modo que buena parte de la etapa final del \emph{pipeline} corrió
  efectivamente en un solo executor, independientemente de $N$.

  \item \textbf{Costo de \emph{exchange} (shuffle) del \emph{join}.}
  Los nodos \texttt{Exchange} muestran 180\,003 registros escritos en
  \emph{shuffle} en el lado de \texttt{tickets}, con un tiempo de
  escritura de hasta 604\,ms en el peor \emph{task}. Al usar
  \texttt{sort-merge join} con difusión desactivada deliberadamente,
  ambas tablas --- incluida \texttt{agentes}, de solo 300 filas --- se
  redistribuyen por red en lugar de aprovechar una unión por difusión.

  \item \textbf{Ordenamiento previo al \emph{merge}.} Los nodos
  \texttt{WholeStageCodegen} de tipo \texttt{Sort} muestran un tiempo de
  ordenamiento de hasta 594\,ms antes de ejecutar el
  \texttt{SortMergeJoin}, etapa inherente a la estrategia de \emph{join}
  cuyo costo no se reduce por agregar executors.

  \item \textbf{Planificación del DAG y evaluación perezosa.} Spark
  construye el plan de ejecución completo antes de materializarlo, y en
  este caso generó 4 \emph{jobs} sucesivos debido a la re-planificación
  adaptativa tras recolectar estadísticas de las etapas iniciales
  \cite{zaharia2016}. Esa coordinación adicional es un costo fijo que
  no escala con los datos ni se beneficia de más executors.
\end{enumerate}

En conjunto, estos cuatro mecanismos --- y en particular el colapso de
particiones por AQE --- muestran que la sobrecarga de T3 no es un
fenómeno difuso atribuible de forma genérica al motor, sino un efecto
localizable y verificable en la propia Spark UI de la ejecución medida.


\section{Propuesta de integración al PFC}


\subsection{Caso de uso}

El PFC ACC (Sistema de Gestión de Soporte Técnico ISP) ya incorpora,
desde su Entrega 3, un \emph{pipeline} de Apache Spark orientado a dos
análisis sobre el histórico de tickets almacenado en
\texttt{ticket\_db} (CockroachDB, clúster de 3 nodos): (i) detección de
\textbf{reincidencia de cliente}, mediante una función de ventana
(\texttt{Window.partitionBy("client\_id").orderBy("timestamp")}) que
calcula, para cada cliente, el tiempo transcurrido desde su incidencia
anterior y marca reincidencia cuando esa diferencia es menor a 30 días;
y (ii) \textbf{clustering} de tickets con K-Means ($k=5$), sobre
\emph{features} de TF-IDF de la descripción textual combinadas con la
zona operativa codificada. Para evitar ambigüedad con la nomenclatura de
PE-U4, estas transformaciones se referencian aquí como
$T_{\mathrm{reincidencia}}$, $T_{\mathrm{flag30d}}$ y
$T_{\mathrm{clustering}}$, distintas de T1--T5 de la práctica grupal.

\subsection{Periodicidad y volumen (declarados como decisión propia)}

El pipeline se ejecutó, hasta ahora, una sola vez de forma manual desde
notebook; no existe una periodicidad definida en el PFC. Se propone,
como decisión de diseño de este informe, una \textbf{ejecución batch
mensual}, alineada con el cierre natural de los reportes gerenciales de
soporte técnico. Para dimensionar esa ejecución se declara el siguiente
supuesto, explícitamente no derivado de datos medidos del PFC: un ISP
mediano opera con un volumen de \textbf{aproximadamente 1\,500 tickets
mensuales} (rango de sensibilidad 500--2\,000, según estacionalidad de
incidencias de red). Este supuesto alimenta el cálculo del umbral de
rentabilidad de la Sección~6.

\subsection{Salida producida y decisión de negocio soportada}

La ejecución mensual del pipeline produce, por cliente: la bandera
$\mathrm{is\_recurring\_30d}$ y el $\mathrm{cluster\_id}$ asignado a
cada ticket, agregados en dos tablas de salida: (a) conteo de clientes
reincidentes por zona y categoría, y (b) distribución de tickets por
clúster y zona. Esta salida soporta una decisión operativa concreta:
priorizar la atención de campo hacia zonas con alta concentración de
clientes reincidentes, en vez de tratar cada ticket como un evento
aislado --- un uso análogo, en espíritu, al enrutamiento automatizado
de tickets hacia el grupo experto correcto que describen Xu et
al.~\cite{xu2018} para sistemas de gestión de TI, aunque aquí el
objetivo no es el enrutamiento en sí sino la identificación proactiva de
patrones de reincidencia antes de que se conviertan en escaladas de SLA.

\subsection{Punto de inserción en la arquitectura en capas}

La arquitectura de ACC es de microservicios con base de datos por
servicio: \texttt{auth-service}, \texttt{ticket-service} y
\texttt{report-service} (Spring Boot/Java) sobre CockroachDB;
\texttt{notification-service} (Node.js) y \texttt{ai-service}
(FastAPI/Python) sobre MongoDB; comunicación asíncrona entre servicios
vía Kafka. El pipeline de Spark se inserta como un \textbf{proceso batch
independiente}, fuera del flujo síncrono de creación/clasificación de
tickets: lee mensualmente por JDBC directamente de \texttt{ticket\_db},
sin pasar por Kafka ni por el \texttt{ai-service}, y escribe sus tablas
de salida en un esquema separado dentro del mismo clúster CockroachDB.
\texttt{report-service} --- que ya expone lectura desacoplada bajo un
patrón CQRS --- añade un \textbf{endpoint nuevo} que consulta esas
tablas de salida y las expone en un panel para el rol
\texttt{Administrador} del frontend (Figura~\ref{fig:arquitectura}).

\begin{figure}[h]
\centering
\includegraphics[width=\linewidth]{figuras/fig_arquitectura_integracion.png}
\caption{Punto de inserción del pipeline de Spark en la arquitectura de
microservicios de ACC. Elaboración propia.}
\label{fig:arquitectura}
\end{figure}

Es deliberado que la integración sea por \textbf{lotes y no en flujo},
pese a que el resto del sistema es orientado a eventos: la detección de
reincidencia requiere una ventana histórica de 30 días por cliente, una
agregación que no tiene sentido recalcular evento por evento en tiempo
real, y que además no puede alimentar de vuelta al \texttt{ai-service}
sin rediseñar su contrato de clasificación síncrona --- una limitación
de integración que se reconoce explícitamente aquí y se retoma como
limitación en la Sección~7.


\section{Dimensionamiento y viabilidad}


El umbral de rentabilidad modela el tiempo secuencial como
$T_{\mathrm{sec}}(n)=\alpha n$ y el distribuido como
$T_{\mathrm{dist}}(n)=\beta+\gamma n/N$, siendo ventajoso el motor
distribuido cuando $T_{\mathrm{dist}}(n) < T_{\mathrm{sec}}(n)$:

\begin{equation}
n^{*} = \frac{\beta}{\alpha - \dfrac{\gamma}{N}},
\qquad \text{válido si } \alpha > \frac{\gamma}{N}.
\label{eq:umbral}
\end{equation}

\textbf{Método de ajuste declarado}: $\hat\alpha$ se estima como
$T_{\mathrm{pandas}}/n$ sobre el punto medido de T3 ($n=600\,000$);
$\hat\beta$ y $\hat\gamma$ se obtienen por regresión lineal de
$T_{\mathrm{dist}}$ contra $1/N$ sobre los tres puntos medidos de T3
($N=1,2,4$), ya que $T_{\mathrm{dist}}=\beta+(\gamma n)\cdot(1/N)$ es
lineal en $1/N$ para $n$ fijo. El resultado es
$\hat\alpha=1.6865\times10^{-5}$ \si{\second}/registro,
$\hat\beta=6.6711$ \si{\second}, $\hat\gamma=1.9280\times10^{-6}$
\si{\second}/registro. Sustituyendo en la ecuación~\eqref{eq:umbral}
con $N=4$:

\begin{equation}
n^{*} = \frac{6.6711}{1.6865\times10^{-5} - \dfrac{1.9280\times10^{-6}}{4}}
       \approx 407{,}206 \text{ registros.}
\end{equation}

Con el volumen mensual declarado en la Sección~5
(\textasciitilde1\,500 tickets/mes), se tardarían aproximadamente 271
meses (\textasciitilde22 años) en acumular ese volumen en una sola
ejecución. \textbf{Recomendación}: Spark no se justifica para el batch
mensual de detección de reincidencia bajo este modelo; se reserva para
el histórico acumulado completo o para un reentrenamiento anual del
clustering, mientras que la detección mensual de reincidencia (ventana
corta, 30 días) se resuelve por debajo del umbral con una consulta SQL
directa sobre CockroachDB (ver alternativa evaluada en la
Sección~7.2). El modelo de servicio en la nube recomendado para las
ejecuciones que sí superen $n^{*}$ es un clúster gestionado de sesión
efímera (por ejemplo, Databricks Job Clusters o AWS EMR Serverless),
justificado por el modelo de costos \emph{pay-as-you-go} que describen
Armbrust et al.~\cite{armbrust2010}: al ejecutarse solo mensualmente o
anualmente, mantener un clúster persistente no se justifica
económicamente frente a la elasticidad de aprovisionamiento bajo
demanda. El principal riesgo identificado es que $\hat\beta$ se estimó
sobre una sesión de Colab, no representativa del costo de arranque real
de un clúster gestionado en producción (ver limitación declarada en la
Sección~7.4).


\section{Postura crítica y alternativas descartadas}

\subsection{¿Es Spark la opción correcta para ACC?}

La respuesta es condicional, no binaria, y se sostiene directamente en
los resultados de las secciones anteriores. A nivel de motor, Spark
demostró ser técnicamente superior a pandas en las cinco
transformaciones medidas en PE-U4 (Tabla~\ref{tab:tiempos}), y el
análisis de Karp-Flatt (Sección~4) permitió diagnosticar con precisión
que su límite de escalado en T3 proviene de sobrecarga de coordinación
--- parcialmente explicada por el colapso adaptativo de particiones
observado en la Spark UI --- y no de una limitación algorítmica del
\emph{join}. En ese plano, Spark es una elección sólida y validada
empíricamente con datos propios, no una moda tecnológica adoptada sin
evidencia.

Sin embargo, el cálculo del umbral de rentabilidad (Sección~6) arroja
$n^*\approx407{,}206$ registros, mientras que el volumen mensual
supuesto para ACC (\textasciitilde1\,500 tickets/mes) tardaría más de
20 años en alcanzarlo. \textbf{Para el caso de uso mensual descrito en
la Sección~5, Spark no está justificado hoy}: el propio modelo
matemático que motiva su adopción, aplicado con honestidad a los
números reales del PFC, recomienda lo contrario a la intuición inicial
del equipo. La postura de este informe es, por tanto, que Spark debe
reservarse para cargas de trabajo que sí crucen ese umbral --- el
histórico acumulado completo, o un reentrenamiento anual del clustering
--- y no para la detección mensual de reincidencia, que se resuelve de
forma más simple y con menor costo operativo sin un motor distribuido.

\subsection{Alternativa 1 descartada: agregación en SQL directamente sobre CockroachDB}

CockroachDB soporta funciones de ventana SQL nativas, equivalentes a la
lógica de Window.partitionBy("client\_id").orderBy("timestamp")
que hoy implementa $T_{\mathrm{reincidencia}}$ en Spark. Para el volumen
mensual de ACC, muy por debajo de $n^*$, una consulta SQL directa sobre
\texttt{ticket\_db} evitaría por completo el costo fijo $\beta$ de
arrancar una sesión Spark, sin sacrificar corrección. Se descarta como
reemplazo definitivo --- no como solución mensual, que sí se recomienda
--- porque no escala si ACC crece varios órdenes de magnitud o si se
agrega el \emph{clustering} sobre texto, que SQL no puede expresar sin
extensiones externas; conviene mantenerla exclusivamente para el
escenario de bajo volumen ya identificado.

\subsection{Alternativa 2 descartada: Dask en lugar de Spark}

Dask ofrece paralelismo de datos en Python puro, sin la sobrecarga de
serialización JVM$\leftrightarrow$Python (Py4J) que Spark introduce
incluso en modo local, y encajaría de forma más homogénea con el stack
Python ya presente en \texttt{ai-service} (FastAPI). Se descarta para
ACC por dos razones técnicas concretas: (i) el equipo ya validó
empíricamente su dominio de Spark en PE-U4, incluida la verificación de
equivalencia numérica pandas$\leftrightarrow$Spark, mientras que migrar
a Dask reiniciaría ese trabajo de validación sin garantía de mejor
resultado a este volumen; y (ii) el conector JDBC de Spark hacia
CockroachDB es más maduro y mejor documentado que el soporte equivalente
en el ecosistema de Dask, lo que reduce el riesgo de integración en el
punto de inserción ya definido en la Sección~5.

\subsection{Limitaciones reconocidas de esta propuesta}

Dos limitaciones deben declararse explícitamente. Primero, el
\emph{clustering} de $T_{\mathrm{clustering}}$ concentra el 81{,}3\,\%
de los registros en un único clúster, un artefacto de la baja
variabilidad léxica de las descripciones generadas por plantilla en el
dataset sintético; su valor de negocio real solo podrá evaluarse con
descripciones de clientes reales, de mayor diversidad textual, en
producción. Segundo, los parámetros $\hat\beta$ y $\hat\gamma$ del
umbral de rentabilidad se estimaron sobre una sesión de Spark en Google
Colab, cuyo costo fijo de arranque puede no ser representativo de un
clúster gestionado en producción (por ejemplo, Databricks o EMR); por
tanto, $n^*=407{,}206$ debe leerse como una estimación de orden de
magnitud, no como un umbral operativo exacto para un despliegue real de
ACC \cite{hillmarty2008}.

\section{Conclusiones}

\begin{enumerate}
  \item El ajuste de la ley de Amdahl sobre T3 ($\hat p=0.4445$,
  Sección~3) predice una curva monótonamente creciente, pero los datos
  medidos muestran una inversión de signo entre $N=2$ y $N=4$
  ($\Delta=+11.49\,\%$ frente a $\Delta=-4.82\,\%$,
  Tabla~\ref{tab:amdahl}): el modelo clásico, que asume sobrecarga nula,
  es insuficiente por sí solo para explicar el comportamiento real de un
  motor distribuido y debe complementarse con un diagnóstico de
  sobrecarga como el de Karp-Flatt.

  \item La fracción serial experimental de T3 crece un 51.95\,\% entre
  $N=2$ y $N=4$ ($e:0.395\rightarrow0.601$, Sección~4), y la evidencia
  directa de la Spark UI --- en particular el colapso adaptativo de
  particiones a una sola partición --- confirma que la causa es
  sobrecarga de coordinación del motor, no un límite algorítmico del
  \emph{join}; esto demuestra que atribuir la pérdida de escalabilidad a
  ``el motor'' de forma genérica es insuficiente cuando existe evidencia
  verificable que apunta a un mecanismo específico.

  \item El umbral de rentabilidad calculado para el PFC ACC
  ($n^*=407{,}206$ registros, Sección~6) es muy superior al volumen
  mensual realista supuesto (\textasciitilde1\,500 tickets/mes,
  Sección~5): la decisión técnicamente honesta, sostenida en el propio
  modelo matemático exigido por esta actividad, es no adoptar Spark
  para el procesamiento mensual de reincidencia, reservándolo para el
  histórico acumulado completo una conclusión que contradice la
  intuición inicial de ``más datos, siempre Spark'' y que solo emerge de
  cuantificar el umbral con datos propios en vez de asumirlo.
\end{enumerate}

\section*{Declaración de uso de inteligencia artificial generativa}

Se utilizó Claude (Anthropic) como herramienta de apoyo en las
siguientes tareas, todas revisadas y validadas por el autor antes de su
inclusión: (i) verificación técnica del repositorio de PE-U4
clonación, comprobación de accesibilidad anónima, cotejo línea por línea
de tiempos crudos.csv/tiempos resumen.csv contra el
commit declarado, y revisión de \texttt{references\_U4.bib} para
confirmar qué referencias bibliográficas eran de aporte propio; (ii) cálculos numéricos ajuste por mínimos cuadrados de la fracción
paralelizable $\hat p$, sustitución numérica de la métrica de
Karp-Flatt, y regresión lineal para estimar $\hat\alpha,\hat\beta,
\hat\gamma$ del umbral de rentabilidad , ejecutados como código
verificable, no estimados a mano; (iii) generación de las figuras
propias (\texttt{fig\_speedup.png} y
\texttt{fig\_arquitectura\_integracion.png}) mediante scripts de
Python, incluida su corrección iterativa tras detectar errores
tipográficos y de dirección lógica en las flechas; y (iv) redacción
asistida de las Secciones~3 a~8, sobre la base de datos, decisiones
arquitectónicas y respuestas de dominio aportadas íntegramente por el
autor.

Se declara explícitamente una divergencia metodológica identificada
durante el proceso: el ajuste de $\hat p=0.4445$ obtenido en este
informe individual, sobre $N=2,4$, difiere del valor $p=0.173423$
reportado en el informe grupal de PE-U4, que empleó una metodología de
ajuste distinta. Esta diferencia se declara para dejar constancia de
que el análisis matemático de este documento es independiente del
grupal, no una reproducción de sus resultados.

\printbibliography

\end{document}
